\documentclass[10pt,twocolumn]{article}

\usepackage[top=0.35in,bottom=0.35in,left=0.4in,right=0.4in]{geometry}
\usepackage[T1]{fontenc}
\usepackage[utf8]{inputenc}
\usepackage{lmodern}
\usepackage{microtype}
\usepackage{titlesec}
\usepackage{titling}

\setlength{\parindent}{0pt}
\setlength{\parskip}{0.28em}
\setlength{\columnsep}{0.22in}
\setlength{\droptitle}{-5.2em}
\pretitle{\begin{center}\Large\bfseries}
\posttitle{\par\end{center}\vspace{-1.15em}}
\preauthor{}
\postauthor{}
\predate{}
\postdate{}
\titlespacing*{\section}{0pt}{0.45em}{0.08em}
\pagestyle{empty}

\title{Glossary}
\author{}
\date{}

\begin{document}

\maketitle

\section{Retrieval Selection}
More than one memory or interpretation can exist at the same time. Retrieval selection is the process by which one of them becomes active enough to guide what a person feels and does.

\section{Presence vs. Selection}
Something can still be real, remembered, or felt, and still not be the thing that guides behavior. Presence means it exists. Selection means it has control in the moment.

\section{Competing Patterns}
Competing patterns are different ways of understanding the same situation that can both be available at once. One may come from what is happening now, while another may come from older memory.

\section{Winner-Take-All}
Winner-take-all means one pattern becomes dominant and suppresses the others. The other patterns do not disappear. They stop guiding what happens.

\section{Inhibitory Control}
Inhibitory control is a filtering system in the brain. It does not store memories. It helps determine which signals are allowed to build and which signals are pushed into the background.

\section{Timing}
Signals do not all arrive at once. Small differences in when they arrive, or when they are allowed through, can decide which pattern takes hold.

\section{Timing-Sensitive Inhibition}
Inhibitory signals arrive in narrow time windows. A signal that arrives slightly earlier or later can block one pattern while allowing another to stabilize.

\section{Oscillatory Control}
The brain operates in rhythms. These rhythms create brief windows where signals can pass or be blocked. Patterns that align with these windows can take hold more easily.

\section{Medial Septum Input}
Medial septum input is a control signal that helps set timing across the memory system. It does not carry the memory itself. It influences which patterns can stabilize.

\section{Medial Entorhinal Cortex}
The medial entorhinal cortex helps shape incoming information before it reaches the final memory-output stage. In this whitepaper, it is treated as a transformation layer where timing and filtering are applied.

\section{CA1 Population Activity}
CA1 population activity is the coordinated activity pattern that reflects which memory or interpretation is active. In this whitepaper, it is treated as the pattern that can guide behavior.

\section{Fear Extinction}
Fear extinction means new experience does not erase an older fear pattern. Both can remain. Depending on context or internal state, the older pattern can return and take control again.

\section{Structure vs. No Structure}
Structure means there are external limits or regular supports around a state-changing input. Regular appointments, limited access, or monitoring can create structure. When that structure changes, the balance of which patterns take hold can also change.

\section{Global Inhibition}
Some substances do not target one specific memory. They broadly change inhibition across the brain. This can shift timing, filtering, and how strongly one pattern settles compared to another.

\section{Settling}
Settling is what happens when the system moves toward one pattern after more than one pattern is possible. Once it settles, that pattern can feel clear and certain, even if it did not feel that way a moment before.

\section{Certainty}
Certainty is the feeling that one interpretation is fully true. In this framework, certainty can come from a pattern stabilizing. It does not always mean the pattern is more accurate than the alternatives.

\section{Non-Erasure}
Non-erasure means an older pattern can remain available after a newer pattern becomes active. The older pattern may stop guiding behavior, but it has not necessarily disappeared.

\section{Behavioral Control}
Behavioral control means a pattern has enough influence to shape what happens next. A pattern can be remembered, missed, or felt without having behavioral control.

\section{Signal vs. Background}
Signal is the pattern the system treats as relevant enough to guide perception or action. Background is material that remains present but does not change the selected interpretation.

\section{State-Dependent Access}
State-dependent access means a memory or interpretation can become easier or harder to retrieve depending on anxiety, arousal, medication context, fatigue, or other internal conditions.

\section{Stabilization Time}
Stabilization time is the time a selected pattern needs before it becomes integrated across moments. A state can feel final before it has had time to settle into a broader pattern.

\section{Mechanism vs. Explanation}
A mechanism describes how a process could work. It does not prove why a person acted in a specific way. This distinction keeps the science from becoming a diagnosis.

\section{Personal Motivation}
Personal motivation is the lived observation that made the scientific question matter. It belongs in the whitepaper as the origin of the question, not as evidence that the mechanism explains one person.

\section{Planes of Inquiry}
Planes of inquiry are the separate levels being held apart: circuit mechanism, personal observation, and therapeutic meaning. Keeping the planes separate helps the whitepaper define a problem without turning it into an accusation or a plea.

\end{document}
